\begin{proposition}[7.2.2]
\label{IV.2.s7.2.2-fr}
Sea $A$ un anillo local noetheriano íntegro; pongamos
$A'=\widehat A$, $X=\operatorname{Spec}(A)$,
$X'=\operatorname{Spec}(A')$ y designemos por $f:X'\to X$ el
morfismo canónico. Sean $a$ el punto cerrado de $X$ y
$j:X\setmin\{a\}\to X$ la inclusión canónica. Sean $\mathscr F$ un
$\mathscr O_X$-módulo coherente y
$\mathscr F'=f^*(\mathscr F)=\mathscr F\otimes_{\mathscr O_X}\mathscr O_{X'}$; entonces son equivalentes las condiciones siguientes:
\begin{enumerate}
\item[\textup{a)}] El $\mathscr O_X$-módulo
$j_*(\mathscr F|X\setmin\{a\})$ es coherente.
\item[\textup{b)}] Para todo $x'\in\operatorname{Ass}(\mathscr F')$ se tiene
$\dim(\overline{\{x'\}})\geqslant2$.
\end{enumerate}
\end{proposition}

\begin{proof}
Sea $a'$ el punto cerrado de $X'$, único punto de la fibra
$f^{-1}(a)$, y sea $j':X'\setmin\{a'\}\to X'$ la inclusión canónica.
Como el morfismo $f$ es fielmente plano y cuasicompacto, decir que
$j_*(\mathscr F|X\setmin\{a\})$ es coherente equivale a decir que
$j'_*(\mathscr F'|X'\setmin\{a'\})$ es coherente (5.9.5).
Por otra parte, como $A'$ es isomorfo al cociente de un anillo local
regular por el teorema de Cohen (\textbf 0, 19.8.8), de (5.11.4) resulta
que la condición \emph{b)} del enunciado equivale a que
$j'_*(\mathscr F'|X'\setmin\{a'\})$ sea coherente. De ahí la
proposición.

En lugar de aplicar (5.11.4) mediante el teorema de Cohen, lo cual
por (5.11.2) recurre implícitamente a los resultados cohomológicos
del capítulo III, también puede utilizarse que $A'$ es
universalmente catenario (5.6.4) y universalmente japonés en virtud
de (7.6.5), cuya demostración no utiliza (7.2.2).
\end{proof}
